% ================================================================
% Yêu cầu thêm vào preamble của report.tex:
%   \usepackage{algorithm}
%   \usepackage{algpseudocode}
%   \usepackage{booktabs}
% ================================================================

% ================================================================
\subsection{Nhận diện tên riêng và địa chỉ (Context Detector)}
% ================================================================

Module \texttt{context\_detector} chịu trách nhiệm phát hiện ba loại thực thể
dựa trên ngữ cảnh xung quanh: tên người (\texttt{NAME}), tên người dùng
(\texttt{USERNAME}) và địa chỉ (\texttt{ADDRESS}).
Khác với \texttt{regex\_detector} vốn chỉ dựa vào cấu trúc cố định của chuỗi
ký tự, module này phân tích ý nghĩa ngữ nghĩa của văn bản để xác định thực thể.
Phần dưới đây tập trung trình bày chi tiết hai thành phần chính: nhận diện tên
người và nhận diện địa chỉ.

% ----------------------------------------------------------------
\subsubsection{Nhận diện tên người}
% ----------------------------------------------------------------

Nhận diện tên người là bài toán khó nhất trong ba loại thực thể vì tên người
không tuân theo một cấu trúc cố định và có thể xuất hiện trong rất nhiều ngữ
cảnh khác nhau. Thuật toán được thiết kế theo hai tầng xử lý: \textit{fast
path} khi có dấu hiệu rõ ràng, và \textit{scoring path} khi cần suy luận từ
ngữ cảnh.

\paragraph{Tiền xử lý --- Tách câu thông minh.}
Trước khi phân tích ngữ cảnh, văn bản đầu vào được tách thành các câu riêng
lẻ. Đây là bước quan trọng vì điểm ngữ cảnh (pronoun, nghề nghiệp, trigger)
chỉ có ý nghĩa trong phạm vi một câu.

Thách thức chính là phân biệt dấu chấm kết thúc câu với dấu chấm trong các
trường hợp khác. Hàm \texttt{\_split\_sentences(text)} xử lý điều này theo quy
tắc: dấu \texttt{.} chỉ được coi là ranh giới câu khi \textit{đồng thời} thỏa
ba điều kiện:

\begin{enumerate}[nosep]
    \item Token kết thúc bằng \texttt{.} không thuộc danh sách viết tắt đã biết
          (\texttt{Dr.}, \texttt{Mr.}, \texttt{St.}, \texttt{Apt.}, \texttt{Inc.}, v.v.).
    \item Phần văn bản ngay sau \texttt{.} không bắt đầu bằng tên miền cấp cao
          (\texttt{com}, \texttt{edu}, \texttt{vn}, \texttt{org}, v.v.) --- nhằm
          tránh split nhầm địa chỉ email và URL.
    \item Sau dấu \texttt{.} là khoảng trắng hoặc hết chuỗi.
\end{enumerate}

Dấu \texttt{!} và \texttt{?} luôn được coi là kết thúc câu vô điều kiện.

\paragraph{Tầng 1 --- Fast path (Trigger mạnh).}
Nếu văn bản chứa một trong các cụm từ trigger, tên được xác định ngay lập tức
mà không cần qua bước chấm điểm. Trigger được chia thành hai nhóm:

\begin{itemize}[nosep]
    \item \textbf{NAME\_TRIGGERS}: cụm từ chỉ thẳng tên chủ thể, gồm
          \texttt{"my name is"}, \texttt{"full name"}, \texttt{"name:"}.
    \item \textbf{\_LABEL\_TRIGGERS}: nhãn hay đứng trước tên trong văn bản có
          cấu trúc, gồm khoảng 25 mẫu như \texttt{"contact information:"},
          \texttt{"author:"}, \texttt{"submitted by:"}, \texttt{"regards,"},
          \texttt{"sincerely,"}, v.v.
\end{itemize}

Sau khi tìm thấy trigger, thuật toán thu thập 1--4 token liên tiếp thỏa:
bắt đầu bằng chữ hoa, không chứa chữ số, không chứa ký tự đặc biệt (ngoại
trừ dấu gạch ngang \texttt{-} vì tên có thể có dạng \texttt{Mary-Jane}),
không phải stopword, và không mang đuôi động từ sau token đầu tiên
(\texttt{-ing}, \texttt{-tion}, \texttt{-ed}, v.v.).

\paragraph{Tầng 2 --- Scoring path (Chấm điểm ngữ cảnh).}
Khi không có trigger, thuật toán chuyển sang scoring path gồm ba bước.

\textbf{Bước 2.1 --- Quét candidate.}
Hàm \texttt{\_extract\_name\_candidates(text)} duyệt toàn bộ văn bản, tìm các
cụm từ viết hoa liên tiếp có độ dài từ 2 đến 4 token. Các điều kiện loại trừ
được áp dụng tại hai mức:

\begin{itemize}[nosep]
    \item \textit{Outer loop} (loại trước khi thu thập): token bắt đầu bằng chữ
          số, token trước là số nguyên thuần túy (dấu hiệu địa chỉ), là title
          prefix (\texttt{Dr.}, \texttt{Mr.}, v.v. --- strip và tiếp tục),
          là \texttt{As} đứng đầu câu, là stopword, hoặc là từ tiêu đề section
          (\texttt{Information}, \texttt{Details}, \texttt{Witness}, v.v.).
    \item \textit{Inner loop} (dừng thu thập giữa chừng): gặp token thường
          (không hoa), chứa ký tự đặc biệt, là chức danh công việc
          (\texttt{Manager}, \texttt{Director}, \texttt{Account}, v.v.), mang
          đuôi động từ sau token đầu tiên, hoặc bắt đầu bằng chữ số.
\end{itemize}

Sau khi thu thập xong, một candidate bị loại nếu token cuối cùng là hậu tố chỉ
tổ chức hoặc địa điểm (\texttt{University}, \texttt{Company}, \texttt{Hospital},
\texttt{Hotel}, v.v.) --- nhóm này gồm khoảng 70 từ chia thành bốn loại: công
ty/doanh nghiệp, tổ chức/cơ quan, trường học, và địa điểm/cơ sở vật chất.

\textbf{Bước 2.2 --- Chấm điểm.}
Hàm \texttt{\_score\_candidate(candidate, sentence, full\_text)} tính điểm cho
từng cặp (candidate, câu). Một candidate xuất hiện qua nhiều câu sẽ lấy điểm
cao nhất trong tất cả các câu đó. Bảng~\ref{tab:scoring} tổng hợp các tín
hiệu và điểm tương ứng.

\begin{table}[H]
    \centering
    \caption{Bảng tín hiệu chấm điểm cho candidate tên người}
    \label{tab:scoring}
    \renewcommand{\arraystretch}{1.3}
    \begin{tabular}{p{8cm} c p{4.5cm}}
        \toprule
        \textbf{Tín hiệu} & \textbf{Điểm} & \textbf{Ghi chú} \\
        \midrule
        Trigger mạnh + candidate đứng ngay sau &
            100, break & Dừng ngay, không so sánh \\
        \texttt{I, [name]} hoặc \texttt{[name], I} &
            100, break & \\
        \texttt{I'm [name]} hoặc \texttt{I am [name]} &
            100, break & \\
        \texttt{I was named [name]} &
            100, break & \\
        Từ \texttt{my name} xuất hiện trong câu &
            +10 & \\
        \texttt{As [name],} đứng đầu câu &
            +5  & \\
        Dấu \texttt{:} đứng ngay trước candidate &
            +5  & Bất kỳ label nào \\
        Appositive: \texttt{[name], <chức danh>,} &
            +3  & Từ sau phẩy là nghề nghiệp \\
        Chữ ký: \texttt{[name]{\textbackslash}n<chức danh>} &
            +4  & Candidate trên dòng riêng \\
        Đại từ ngôi 1 (\texttt{I}, \texttt{me}, \texttt{my}) trong câu &
            +2  & Ưu tiên hơn ngôi 3 \\
        Đại từ ngôi 3 (\texttt{he}, \texttt{she}, \texttt{his}, \texttt{her}) &
            +1  & Chỉ tính khi không có ngôi 1 \\
        Nghề nghiệp xuất hiện trong cùng câu &
            +2  & Xem phần phát hiện nghề nghiệp \\
        Tên lặp lại trong toàn đoạn (partial match) &
            +1/lần & Tối đa +5 \\
        Title prefix (\texttt{Dr.}, \texttt{Mr.}, v.v.) trước tên &
            +1  & \\
        \bottomrule
    \end{tabular}
\end{table}

Lưu ý rằng đại từ không cộng dồn nếu xuất hiện nhiều lần trong cùng một câu.
Nếu câu đồng thời có đại từ ngôi 1 và ngôi 3 thì chỉ tính điểm cho ngôi 1.

\textbf{Phát hiện nghề nghiệp.}
Hàm \texttt{\_has\_job\_in\_sentence(sentence)} trả về \texttt{True} nếu câu
chứa tín hiệu nghề nghiệp theo một trong hai cách:
(1)~từ nằm trong danh sách cứng khoảng 35 nghề (\texttt{doctor}, \texttt{lawyer},
\texttt{developer}, \texttt{CEO}, v.v.) bất kể vị trí trong câu;
(2)~từ mang đuôi nghề nghiệp (\texttt{-er}, \texttt{-ist}, \texttt{-or},
\texttt{-ician}, v.v.) \textit{và} đứng ngay sau một role marker như \texttt{a},
\texttt{an}, \texttt{is a}, \texttt{works as}. Điều kiện role marker giúp tránh
false positive như \texttt{underwater} hay \texttt{after}.

\textbf{Bước 2.3 --- Chọn winner.}
Sau khi chấm điểm toàn bộ candidate, winner được chọn theo quy tắc:
nếu có candidate đạt điểm 100 (definitive) thì trả về ngay;
nếu chỉ có một candidate thì trả về candidate đó dù điểm bằng 0;
nếu nhiều candidate thì chọn điểm cao nhất, tie-breaking theo vị trí
xuất hiện sớm hơn trong văn bản.

\begin{algorithm}[H]
\caption{Thuật toán nhận diện tên người}
\label{alg:detect_name}
\begin{algorithmic}[1]
\Require Văn bản đầu vào $T$
\Ensure Danh sách thực thể tên $E$

\State $E \leftarrow \emptyset$
\State $lower \leftarrow \text{lowercase}(T)$

\Statex \hspace{-\algorithmicindent}\textbf{// Tầng 1: Fast path}
\For{mỗi trigger $tr$ trong \textsc{NameTriggers} $\cup$ \textsc{LabelTriggers}}
    \For{mỗi vị trí $pos$ của $tr$ trong $lower$}
        \State Thu thập 1--4 token hoa liên tiếp sau $pos$ vào $tokens$
        \If{$tokens \neq \emptyset$}
            \State Thêm $\text{join}(tokens)$ vào $E$ với điểm 100
        \EndIf
    \EndFor
\EndFor
\If{$E \neq \emptyset$} \Return $E$ \EndIf

\Statex \hspace{-\algorithmicindent}\textbf{// Tầng 2: Scoring path}
\State $sentences \leftarrow \textsc{SplitSentences}(T)$
\State $candidates \leftarrow \textsc{ExtractCandidates}(T)$
\If{$candidates = \emptyset$} \Return $\emptyset$ \EndIf

\State $bestScores \leftarrow \{\}$ \quad \textit{// candidate $\to$ (score, start, end)}

\For{mỗi candidate $c$ trong $candidates$}
    \State $best \leftarrow 0$
    \For{mỗi câu $s$ trong $sentences$}
        \If{không có phần nào của $c$ trong $s$} \textbf{continue} \EndIf
        \State $(score, def) \leftarrow \textsc{ScoreCandidate}(c, s, T)$
        \If{$def = \text{True}$}
            \State \Return $\{c\}$ \quad \textit{// Definitive, dừng ngay}
        \EndIf
        \State $best \leftarrow \max(best, score)$
    \EndFor
    \State Cập nhật $bestScores[c]$ nếu $best$ cao hơn giá trị hiện tại
\EndFor

\State $winner \leftarrow \arg\max_{c}\ bestScores[c]$ \quad \textit{// Tie: chọn cái xuất hiện sớm hơn}
\State \Return $\{winner\}$
\end{algorithmic}
\end{algorithm}

% ----------------------------------------------------------------
\subsubsection{Nhận diện địa chỉ}
% ----------------------------------------------------------------

Địa chỉ trong tập dữ liệu tuân theo cấu trúc điển hình của địa chỉ Mỹ gồm
số nhà, tên đường, loại đường, và hậu tố hướng hoặc đơn vị căn hộ tùy chọn.
Module nhận diện địa chỉ sử dụng hai cơ chế bổ sung cho nhau.

\paragraph{Cơ chế 1 --- Regex pattern.}
Biểu thức chính quy \texttt{\_ADDRESS\_RE} mô tả cấu trúc địa chỉ theo bốn
thành phần nối tiếp:

\begin{enumerate}[nosep]
    \item \textbf{Số nhà}: \verb|\b(\d{1,5})(?!\d)\s+| ---
          từ 1 đến 5 chữ số, kèm negative lookahead \verb|(?!\d)| để đảm
          bảo đây không phải là một phần của số dài hơn (loại trừ số điện
          thoại, mã bưu điện dài).

    \item \textbf{Tên đường}: \verb|((?:[A-Z0-9][A-Za-z0-9]*\s+)+?)| ---
          một hoặc nhiều từ bắt đầu bằng chữ hoa hoặc chữ số, sử dụng
          non-greedy để nhường ưu tiên cho thành phần loại đường phía sau.

    \item \textbf{Loại đường}: danh sách khoảng 30 từ được sắp xếp từ dài
          đến ngắn để tránh match sớm, bao gồm dạng đầy đủ
          (\texttt{Boulevard}, \texttt{Terrace}, \texttt{Avenue},
          \texttt{Street}, \texttt{Drive}, \texttt{Road}, v.v.) và dạng viết
          tắt (\texttt{Blvd.}, \texttt{Ave.}, \texttt{Rd.}, \texttt{St.}, v.v.).

    \item \textbf{Hướng và đơn vị} (tùy chọn):
          hướng gồm tám giá trị \texttt{North}, \texttt{South}, \texttt{East},
          \texttt{West} và các tổ hợp; đơn vị gồm \texttt{Suite},
          \texttt{Apt.}, \texttt{Apartment}, \texttt{Unit} kèm số phòng.
\end{enumerate}

Trước khi kiểm tra kết quả regex, tất cả span email trong văn bản được xác định
trước; bất kỳ span địa chỉ nào nằm hoàn toàn bên trong span email sẽ bị bỏ qua
để tránh nhầm lẫn phần tên miền của địa chỉ email với địa chỉ nhà.

\paragraph{Cơ chế 2 --- Trigger keyword (fallback).}
Đối với địa chỉ không có loại đường tiêu chuẩn, module sử dụng danh sách
trigger bao gồm \texttt{street}, \texttt{road}, \texttt{avenue},
\texttt{district}, \texttt{city}, \texttt{st.}, \texttt{rd.}, \texttt{ave.}.
Khi tìm thấy trigger, thuật toán lấy clause xung quanh (từ dấu câu trước đến
dấu câu sau), loại bỏ noise words đầu clause (\texttt{at}, \texttt{is},
\texttt{the}, \texttt{located}, \texttt{live}, \texttt{reside}, v.v.), và kiểm
tra thêm: số nhà ở đầu clause không được vượt quá 5 chữ số, và span không được
trùng lặp với kết quả từ cơ chế regex.

\begin{algorithm}[H]
\caption{Thuật toán nhận diện địa chỉ}
\label{alg:detect_address}
\begin{algorithmic}[1]
\Require Văn bản đầu vào $T$
\Ensure Danh sách thực thể địa chỉ $E$

\State $E \leftarrow \emptyset$
\State $emailSpans \leftarrow \textsc{FindEmailSpans}(T)$

\Statex \hspace{-\algorithmicindent}\textbf{// Cơ chế 1: Regex pattern}
\For{mỗi match $m$ của \textsc{AddressRE} trong $T$}
    \If{$m$ nằm trong một email span} \textbf{continue} \EndIf
    \State Thêm span của $m$ vào $E$
\EndFor

\Statex \hspace{-\algorithmicindent}\textbf{// Cơ chế 2: Trigger keyword fallback}
\For{mỗi trigger $tr$ trong \textsc{AddressTriggers}}
    \For{mỗi vị trí $pos$ của $tr$ trong $T$}
        \State $clause \leftarrow \textsc{ExtractClause}(T, pos)$
        \State $addr \leftarrow \textsc{StripLeadingNoise}(clause)$
        \If{số nhà đầu $addr$ có nhiều hơn 5 chữ số} \textbf{continue} \EndIf
        \If{$addr$ nằm trong email span} \textbf{continue} \EndIf
        \If{$addr$ overlap với span đã có trong $E$} \textbf{continue} \EndIf
        \State Thêm span của $addr$ vào $E$
    \EndFor
\EndFor

\State \Return $\textsc{Deduplicate}(E)$
\end{algorithmic}
\end{algorithm}

% ================================================================
\subsection{Hợp nhất thực thể (Merger)}
% ================================================================

Sau khi \texttt{regex\_detector} và \texttt{context\_detector} chạy độc lập,
mỗi module trả về một danh sách thực thể riêng. Hai danh sách này được đưa vào
module \texttt{merger} để hợp nhất thành một danh sách duy nhất, sạch và
không chồng lấn.

\subsubsection{Vấn đề overlap}

Hai thực thể được coi là \textit{chồng lấn} (overlap) nếu khoảng ký tự của
chúng giao nhau, tức là $[s_1, e_1)$ và $[s_2, e_2)$ thỏa $s_1 < e_2$ và
$e_1 > s_2$. Overlap xảy ra trong nhiều tình huống thực tế:

\begin{itemize}[nosep]
    \item Địa chỉ email \texttt{kazuosun@hotmail.net} bị regex\_detector bắt
          đúng là \texttt{EMAIL}, nhưng context\_detector có thể nhận nhầm
          phần \texttt{hotmail.net} là URL hoặc nhận phần trước \texttt{@} là
          \texttt{NAME}.
    \item Số điện thoại \texttt{+27 68 670 7513} có thể bị nhận vừa là
          \texttt{PHONE} vừa overlap với một span \texttt{ADDRESS} nếu đứng
          sau một địa chỉ.
    \item Tên người \texttt{Baha Hoffman} xuất hiện ngay trước số điện thoại,
          context\_detector có thể kéo dài span sang phần số.
\end{itemize}

\subsubsection{Hệ thống ưu tiên}

Để giải quyết overlap một cách nhất quán, merger định nghĩa một thứ tự ưu tiên
tuyến tính cho sáu loại thực thể:

\begin{center}
\texttt{EMAIL} $>$ \texttt{URL} $>$ \texttt{PHONE} $>$
\texttt{ADDRESS} $>$ \texttt{NAME} $>$ \texttt{USERNAME}
\end{center}

Thứ tự này phản ánh độ tin cậy của từng detector: \texttt{EMAIL}, \texttt{URL}
và \texttt{PHONE} được phát hiện bằng regex có độ chính xác cao và cấu trúc
rõ ràng nên được ưu tiên hơn; \texttt{ADDRESS}, \texttt{NAME},
\texttt{USERNAME} dựa trên ngữ cảnh nên được xếp sau.

\subsubsection{Thuật toán Greedy Priority-First}

Hàm \texttt{resolve\_overlaps} thực hiện giải quyết overlap theo chiến lược
tham lam ưu tiên (greedy priority-first):

\begin{enumerate}[nosep]
    \item Gộp tất cả thực thể từ mọi detector thành một danh sách duy nhất.
    \item Sắp xếp danh sách theo ba tiêu chí lần lượt: (a)~thứ hạng ưu tiên
          tăng dần (loại ưu tiên cao xếp trước), (b)~chỉ số \texttt{start}
          tăng dần, (c)~độ dài span giảm dần (span dài hơn xếp trước khi cùng
          điểm xuất phát).
    \item Duyệt tuần tự qua danh sách đã sắp xếp. Với mỗi thực thể, giữ lại
          nếu span của nó không chồng lấn với bất kỳ thực thể đã giữ nào;
          bỏ qua nếu ngược lại.
    \item Sắp xếp lại danh sách kết quả theo \texttt{start} để đảm bảo thứ tự
          xuất hiện trong văn bản.
\end{enumerate}

\begin{algorithm}[H]
\caption{Thuật toán hợp nhất và giải quyết overlap (Greedy Priority-First)}
\label{alg:merger}
\begin{algorithmic}[1]
\Require Các nhóm thực thể $G_1, G_2, \ldots, G_k$ từ các detector
\Ensure Danh sách thực thể không overlap, sắp theo vị trí

\State $all \leftarrow G_1 \cup G_2 \cup \cdots \cup G_k$

\State Sắp xếp $all$ theo $\bigl(\text{priority}(type),\ start,\ -(end - start)\bigr)$ tăng dần

\State $kept \leftarrow \emptyset$
\State $occupied \leftarrow \emptyset$ \quad \textit{// Tập các span $(s, e)$ đã giữ}

\For{mỗi thực thể $ent$ trong $all$}
    \State $s \leftarrow ent.start$,\quad $e \leftarrow ent.end$
    \If{$\nexists\, (os, oe) \in occupied$ sao cho $s < oe$ và $e > os$}
        \State $kept \leftarrow kept \cup \{ent\}$
        \State $occupied \leftarrow occupied \cup \{(s, e)\}$
    \EndIf
\EndFor

\State Sắp xếp $kept$ theo $start$ tăng dần
\State \Return $kept$
\end{algorithmic}
\end{algorithm}

\subsubsection{Minh họa trường hợp overlap}

Xét đoạn văn bản sau được trích từ tập dữ liệu:

\begin{center}
\textit{``please do not hesitate to contact Baha Hoffman at +27 68 670 7513
or via email at \texttt{bahahoffman@yahoo.net}.''}
\end{center}

Hai detector tạo ra các thực thể như trong Bảng~\ref{tab:overlap_before}.

\begin{table}[H]
    \centering
    \caption{Danh sách thực thể trước khi hợp nhất}
    \label{tab:overlap_before}
    \renewcommand{\arraystretch}{1.3}
    \begin{tabular}{l l l c}
        \toprule
        \textbf{Detector} & \textbf{Loại} & \textbf{Text} & \textbf{Ưu tiên} \\
        \midrule
        regex\_detector & \texttt{PHONE} & \texttt{+27 68 670 7513} & 3 \\
        regex\_detector & \texttt{EMAIL} & \texttt{bahahoffman@yahoo.net} & 1 \\
        context\_detector & \texttt{NAME}  & \texttt{Baha Hoffman} & 5 \\
        context\_detector & \texttt{NAME}  & \texttt{Baha Hoffman at} & 5 \\
        \bottomrule
    \end{tabular}
\end{table}

Sau khi sắp xếp theo ưu tiên và áp dụng thuật toán greedy, merger loại bỏ
\texttt{Baha Hoffman at} vì overlap với \texttt{Baha Hoffman} đã được giữ
trước đó (cùng loại \texttt{NAME} nhưng span ngắn hơn xếp trước).
Kết quả cuối cùng được trình bày trong Bảng~\ref{tab:overlap_after}.

\begin{table}[H]
    \centering
    \caption{Danh sách thực thể sau khi hợp nhất}
    \label{tab:overlap_after}
    \renewcommand{\arraystretch}{1.3}
    \begin{tabular}{l l l}
        \toprule
        \textbf{Loại} & \textbf{Text} & \textbf{Vị trí} \\
        \midrule
        \texttt{NAME}  & \texttt{Baha Hoffman}            & start=35, end=47  \\
        \texttt{PHONE} & \texttt{+27 68 670 7513}          & start=51, end=66  \\
        \texttt{EMAIL} & \texttt{bahahoffman@yahoo.net}    & start=79, end=103 \\
        \bottomrule
    \end{tabular}
\end{table}

Kết quả này sau đó được đưa vào module \texttt{anonymizer} để thay thế các
span PII bằng placeholder tương ứng, tạo ra văn bản đã được ẩn danh hóa.